%\smallskip
%\noindent\textbf{Existence of Nash equilibria.}
%Continuous bidding makes the Markov bidding game have discontinuous payoffs at bid ties, since an arbitrarily small increase in bid by any of the tied policies will change the winning policy.
%Simon and Zame~\citep{simon1990discontinuous} prove the existence of Nash equilibria for games with discontinuous payoffs by allowing the tie-breaking mechanism (formalized as ``endogenous payoff sharing rule''  at discontinuity points) to be computed as part of the equilibrium object.
%This does not directly imply that a Nash equilibrium would exist for any arbitrary tie-breaking mechanism fixed apriori, which is the case for us. 
%Nevertheless, we chose uniform tie-breaking due to its simplicity, and since the event that a tie happens has Lebesgue measure zero, in practice, 
%
%in the continuous all-pay equilibrium characterized below, bid ties occur with probability zero, so the particular tie-breaking rule is payoff-irrelevant for that equilibrium.
%The winner-pays result below should instead be read conditionally on the existence or selection of a pure equilibrium.

\smallskip
\noindent\textbf{Equilibrium continuation values.}
We consider local policies \(\pi_i=(\pi_i^a,\pi_i^b)\) that form a Nash equilibrium of the Markov bidding game $\mathcal{G}^{\mathcal{M}}_{\tau,\beta,\rho,\mathsf{P}}$.
For simplicity, assume $H$ is a multiple of $\tau$ (otherwise, extend the horizon to the next multiple of $\tau$).
The policies bid at \textit{bidding rounds} $t\in \{0,\tau,2\tau,\ldots,H-\tau\}$.
Given a policy index $i\in [1;m]$, time $t\leq H$, and state $s\in \hat S$ of $\mathcal{G}^{\mathcal{M}}_{\tau,\beta,\rho,\mathsf{P}}$, we write $V^*_{i,t}$ to denote the value policy $i$ achieves in the rest of the game starting at time $t$ and at state $s$, when all the policies follow a Nash equilibrium.
Since the horizon $H$ is finite, we can define $V^*_{i,t}$ using backward induction over $t$ starting from $t=H$.
Set \(V^*_{i,H}(s)=0\) for every \(i\) and \(s\).
Suppose the game is at state \(s\) and bidding round $t\in \{0,\tau,2\tau,\ldots,H-\tau\}$.

The \textit{winning continuation value} of a policy $i$ is the optimal value it can achieve upon winning bidding:
\begin{equation}
	\widetilde V^{\win}_{i,t}(s)
	\coloneqq
	\max_{\pi_i^a}
	\E^{\pi_i^a}\left[
		\sum_{q=0}^{\tau-1} r_i(s^{t+q},a^{t+q})
		+V^*_{i,t+\tau}(s^{t+\tau})
	\right],
\end{equation}
where the action sequence is generated by policy \(i\) during the time window $[t;t+\tau]$.
Similarly, the \textit{losing continuation value} is the least value it achieves upon losing the bidding:
\begin{equation}\label{eq:losing continuation value}
	\widetilde V^{\lose}_{i,t}(s)
	\coloneqq
	\min_{j\neq i}
	\E^{\pi_j^a}\left[
		\sum_{q=0}^{\tau-1} r_i(s^{t+q},a^{t+q})
		+V^*_{i,t+\tau}(s^{t+\tau})
	\right],
\end{equation}
where the action sequence during $[t;t+\tau]$ is generated by a different policy \(j\) using an action policy $\pi_j^a$ that maximizes $j$'s own winning continuation value.
The \textit{control gain} of policy \(i\) at \((s,t)\) is
\begin{equation}
	G_{i,t}(s)
	\coloneqq
	\widetilde V^{\win}_{i,t}(s)-\widetilde V^{\lose}_{i,t}(s).
\end{equation}
This quantity measures the advantage of winning control for the next bidding interval rather than losing it.
Given bids \(b=(b_1,\ldots,b_m)\), let
$W(b)\coloneqq \{k\in [1;m]\mid b_k=\max_{\ell\in[1;m]}b_\ell\}$ 
be the set of highest bidders, and let \(p_i(b)=\mathbf{1}[i\in W(b)]/\abs{W(b)}\) be the probability that \(i\) wins after uniform tie-breaking.
For a penalty model \(\mathsf{P}\in\{\poorman,\allpay\}\), the \textit{one-round bidding payoff} is
\begin{equation}
	U^{\mathsf{P}}_{i,t}(s,b)
	=
	\begin{cases}
		\widetilde V^{\lose}_{i,t}(s)+p_i(b)\left(G_{i,t}(s)-\rho b_i\right) & \mathsf{P}=\poorman,\\
		\widetilde V^{\lose}_{i,t}(s)+p_i(b)G_{i,t}(s)-\rho b_i & \mathsf{P}=\allpay.
	\end{cases}
\end{equation}
In the \poorman case, the bid cost is paid only in expectation over winning; in the \allpay case, the bid cost is paid regardless of the winner.
Effectively, in the Markov game, the goal of policy $i$ boils down to maximizing the one-round bidding payoff $U^{\mathsf{P}}_{i,t}(s,b)$ at each state $s$ and bidding round $t$. 

A mixed bidding policy for player \(i\) at \((s,t)\) is the bid component \(\pi^b_{i,t}(\cdot\mid s)\in\dist([0,\beta])\), and a mixed bidding profile is \(\pi^b_t(\cdot\mid s)=(\pi^b_{1,t}(\cdot\mid s),\ldots,\pi^b_{m,t}(\cdot\mid s))\).
We write \(b\sim\pi^b_t(\cdot\mid s)\) when bids are sampled independently as \(b_i\sim\pi^b_{i,t}(\cdot\mid s)\).
A mixed bidding profile \(\pi^{b*}_t(\cdot\mid s)\) is a Nash equilibrium of the one-round bidding game at \((s,t)\) if, for every player \(i\) and every alternative bid distribution $\pi^b_{i,t}$, where \(\pi^b_{i,t}(\cdot\mid s)\in\dist([0;\beta])\) for $\mathsf{P}=\poorman$ and $\mathsf{P}=\allpay$,
\begin{equation}
	\E_{b\sim\pi^{b*}_t(\cdot\mid s)}\!\left[U^{\mathsf{P}}_{i,t}(s,b)\right]
	\geq
	\E_{b_i\sim\pi^b_{i,t}(\cdot\mid s),\; b_{-i}\sim\pi^{b*}_{-i,t}(\cdot\mid s)}\!\left[U^{\mathsf{P}}_{i,t}(s,b_i,b_{-i})\right].
\end{equation}
For the \poorman setting, the existence of the Nash equilibrium is guaranteed due to the finiteness of the action space $\hat A$~\citep{nash1951noncooperative}, and for the \allpay setting, the same follows from known results in the literature~\citep{baye1996all}.
The equilibrium continuation value is then
\begin{equation}
	V^*_{i,t}(s)
	=
	\E_{b\sim \pi^{b*}_t(\cdot\mid s)}\left[U^{\mathsf{P}}_{i,t}(s,b)\right],
\end{equation}
which completes the backward-induction definition.

\smallskip
\noindent\textbf{Theoretical analysis of the bidding schemes.}
We will now prove that, if all policies follow the Nash equilibrium and if some mild assumptions hold true, then the policy with the highest control gain is guaranteed to get the control at each state and time.
Consequently, we obtain a greedy solution to maximize the global reward (the sum of local cumulative rewards).
Further, from Eq.~\eqref{eq:losing continuation value} it follows that the control gains are actually environment-aware, i.e., if the goal of policy $i$ is aligned with every other policy (like for the cat feeder, all cats are crowded in one corner), then the control gain of $i$ would be automatically low.
We state our two results here but defer the proofs to~\Cref{app:poorman-control-gain-proof,app:allpay-control-gain-proof}.


\begin{theorem}[\poorman favors high control gain]\label{thm:poorman-favors-high-control-gain}
	In the \poorman setting, consider a single bidding round $t$ at the agent state $s\in S$.
  	Suppose the parameters $\rho$ and $\beta$ are so chosen that $\rho\beta > \max_{j\in [1;m]}G_{j,t}(s)$.
	Let $i^*$ be the policy with the largest control gain, such that the following margin condition is fulfilled:
	\begin{align}\label{eq:all pay:gross gain gap}
		G_{i^*,t}(s) - \max_{j\neq i^*}G_{j,t}(s) > \rho.
	\end{align}
	Then in any pure Nash equilibrium, policy $i^*$ wins control.
\end{theorem}

\begin{theorem}[\allpay favors high control gain]\label{thm:allpay-favors-high-control-gain}
	In the \allpay setting with $m\geq 3$, consider a single bidding round $t$ at the agent state $s\in S$.
	Let $i^*$ and $j^*$ be policies for which the following gross gain order holds:
	\begin{align}\label{eqn:allpay urgency gap}
		G_{i^*,t} > G_{j^*,t} > \max_{k\notin \{i^*,j^*\}}G_{k,t}.
	\end{align}
	Then there is no Nash equilibrium with pure strategies, although a unique Nash equilibrium exists with mixed strategies.
	In this Nash equilibrium, policy $i^*$ wins control with probability $1-\sfrac{G_{j^*,t}}{2G_{i^*,t}}$.	
\end{theorem}

While these results do not automatically imply equivalence to the centralized optimal solution, we obtain an environment-aware, greedy approximation.
This gives a scheduling perspective on the auction: each bidding round picks the objective to receive control based on its urgency.
In Appendix~\ref{app:lstf-analysis}, we prove that in deterministic environments with stationary targets and varying deadlines, a Nash equilibrium in our auction mechanism always selects critically urgent goals, if any.
This matches the least slack time first (LSTF) rule, a widely used scheduling heuristic~\citep{brown2020optimal}.


%
%\smallskip
%\noindent\textbf{Theoretical analysis of the bidding schemes.}
%The first result says that the winner-pays mechanism can implement the largest-gain policy in arbitrarily good pure approximate equilibria.
%Recall that a pure bid profile is a pure \(\varepsilon\)-Nash equilibrium if no policy can improve its payoff by more than \(\varepsilon\) through a unilateral bid deviation.
%
%\begin{theorem}[Winner-pays approximate implementation]
%	In the \poorman setting, consider a single bidding round at state \(s\) and time \(t\).
%	Let \(i^*\) be a policy with largest control gain, let \(j^*\) be a policy with second-largest control gain, and suppose
%	\begin{equation}
%		G_{i^*,t}(s)>G_{j^*,t}(s).
%	\end{equation}
%	For any \(\varepsilon>0\) such that \((G_{j^*,t}(s)+\varepsilon)/\rho\leq \beta\) and \(G_{i^*,t}(s)>G_{j^*,t}(s)+\varepsilon\), there is a pure \(\varepsilon\)-Nash equilibrium in which policy \(i^*\) wins control with probability one.
%\end{theorem}
%
%\begin{proof}
%	Consider the bid profile \(b_{j^*}=G_{j^*,t}(s)/\rho\), \(b_{i^*}=(G_{j^*,t}(s)+\varepsilon)/\rho\), and \(b_k\leq b_{j^*}\) for every \(k\notin\{i^*,j^*\}\).
%	Policy \(i^*\) wins uniquely.
%	No losing policy can gain by deviating to another losing bid.
%	If \(j^*\) deviates to a winning bid, it must bid at least \(b_{i^*}\), giving payoff at most
%	\begin{equation}
%		\widetilde V^{\lose}_{j^*,t}(s)+G_{j^*,t}(s)-\rho b_{i^*}
%		=
%		\widetilde V^{\lose}_{j^*,t}(s)-\varepsilon,
%	\end{equation}
%	which is worse than losing.
%	Every other losing policy has control gain at most \(G_{j^*,t}(s)\), so the same argument applies.
%	
%	Policy \(i^*\) cannot gain by raising its bid.
%	Its best deviation is to lower its bid while still winning, which can improve its payoff by at most \(\varepsilon\), since it must bid above \(b_{j^*}\) to remain the unique winner.
%	Deviating to a losing bid gives payoff \(\widetilde V^{\lose}_{i^*,t}(s)\), while the proposed bid gives
%	\begin{equation}
%		\widetilde V^{\lose}_{i^*,t}(s)+G_{i^*,t}(s)-G_{j^*,t}(s)-\varepsilon,
%	\end{equation}
%	which is strictly larger by assumption.
%	Thus no policy can improve by more than \(\varepsilon\), so the profile is a pure \(\varepsilon\)-Nash equilibrium.
%\end{proof}
%
%The all-pay mechanism admits a standard mixed-equilibrium construction in the continuous complete-information setting.
%
%\begin{theorem}[All-pay concentration]
%	In the \allpay setting, consider a single bidding round at state \(s\) and time \(t\), with continuous bid space \([0,\beta]\).
%	Let \(i^*\) and \(j^*\) be policies satisfying
%	\begin{equation}
%		G_{i^*,t}(s)
%		>
%		G_{j^*,t}(s)
%		>
%		\max_{k\notin\{i^*,j^*\}}G_{k,t}(s),
%	\end{equation}
%	and suppose \(\rho\beta\geq G_{j^*,t}(s)\).
%	Then the continuous all-pay bidding game has no pure Nash equilibrium, and there is a mixed Nash equilibrium given by the following bid distributions:
%	\begin{itemize}
%		\item policy \(i^*\) samples uniformly from \([0,G_{j^*,t}(s)/\rho]\);
%		\item policy \(j^*\) places an atom of mass
%		\begin{equation}
%			\frac{G_{i^*,t}(s)-G_{j^*,t}(s)}
%			{G_{i^*,t}(s)}
%		\end{equation}
%		at \(0\), and otherwise has density \(\rho/G_{i^*,t}(s)\) on \((0,G_{j^*,t}(s)/\rho]\);
%		\item every other policy bids \(0\) with probability one.
%	\end{itemize}
%	In this equilibrium, policy \(i^*\) wins control with probability
%	\begin{equation}
%		1-\frac{G_{j^*,t}(s)}
%		{2G_{i^*,t}(s)}.
%	\end{equation}
%\end{theorem}
%
%\begin{proof}
%	This is a standard complete-information all-pay equilibrium~\citep{baye1996all}, applied to valuations \(G_{i,t}(s)\) and bid cost \(\rho b_i\).
%	The bid-cap assumption ensures that the equilibrium support lies inside \([0,\beta]\).
%	The stated winning probability follows by integrating the probability that \(j^*\)'s bid is below the bid of \(i^*\):
%	\begin{equation}
%	\begin{aligned}
%		P[i^*\text{ wins}]
%		&=
%		\int_0^{G_{j^*,t}(s)/\rho}
%		\left(
%			\frac{G_{i^*,t}(s)-G_{j^*,t}(s)}
%			{G_{i^*,t}(s)}
%			+
%			\frac{\rho b}{G_{i^*,t}(s)}
%		\right)
%		\frac{\rho}{G_{j^*,t}(s)}
%		db\\
%		&=
%		1-\frac{G_{j^*,t}(s)}
%		{2G_{i^*,t}(s)}.
%	\end{aligned}
%	\end{equation}
%\end{proof}
